\PrelimChapter{Resumen}{\topBanner}
\vspace*{0.7cm}

El presente trabajo desarrolla la implementación de un servicio de directorio para la infraestructura tecnológica del grupo de investigación en Redes, Información y Distribución (GRID) de la Universidad del Quindío. La propuesta surge ante la necesidad de centralizar la administración de usuarios, credenciales y accesos, dentro de un entorno académico en el que las plataformas tecnológicas del grupo gestionan las identidades de forma independiente, lo que incrementa los procesos manuales de administración y limita la aplicación uniforme de los controles de seguridad.

Metodológicamente, el trabajo comprende: (a) la identificación de las necesidades, problemas y oportunidades (NPO) del grupo GRID con relación a un servicio de directorio, (b) la caracterización de alternativas de software libre y código abierto orientadas a la implementación del servicio, (c) la aplicación del proceso de Análisis y Resolución de Decisiones (DAR) del modelo CMMI para la selección de la tecnología, (d) el diseño del servicio de directorio mediante el lenguaje de modelado ArchiMate y (e) la implementación y validación de un prototipo funcional.

Como resultado, se seleccionó una plataforma integral de gestión de identidades y accesos que articula un servicio de directorio basado en el protocolo ligero de acceso a directorios (LDAP), autenticación Kerberos, gestión de certificados digitales y políticas centralizadas de control de acceso, desplegada bajo un esquema de alta disponibilidad con balanceo de carga entre un servidor principal y un servidor réplica. Finalmente, se implementó un prototipo funcional integrado mediante LDAP con los servicios OpenVPN, OpenProject y Xen Orchestra, validando la viabilidad técnica y funcional de la propuesta para el fortalecimiento de la seguridad y la administración eficiente de los activos de información del grupo GRID.

\PrelimChapter{Abstract}{\topBanner}
\vspace*{0.7cm}

This work develops the implementation of a directory service for the technological infrastructure of the Research Group on Networks, Information and Distribution (GRID) at the University of Quindío. The proposal arises from the need to centralize the administration of users, credentials, and access within an academic environment where the group's technological platforms manage identities independently, which increases manual administration processes and limits the uniform application of security controls.

Methodologically, the work comprises: (a) the identification of the needs, problems, and opportunities (NPO) of the GRID group regarding a directory service, (b) the characterization of free and open source alternatives oriented toward the implementation of the service, (c) the application of the Decision Analysis and Resolution (DAR) process from the CMMI model for technology selection, (d) the design of the directory service using the ArchiMate modeling language, and (e) the implementation and validation of a functional prototype.

As a result, an integrated identity and access management platform was selected, articulating a directory service based on the Lightweight Directory Access Protocol (LDAP), Kerberos authentication, digital certificate management, and centralized access control policies, deployed under a high availability scheme with load balancing between a primary server and a replica server. Finally, a functional prototype integrated through LDAP with the OpenVPN, OpenProject, and Xen Orchestra services was implemented, validating the technical and functional viability of the proposal for strengthening the security and efficient administration of the information assets of the GRID group.
